(i) $\omega^2 = (-1)^3(-1)^1 = +1$ and $p - q = 1$: type C by Fact \claimref{rule2.imported-classification-centre-type-c}; the corner is $\Cl^0(2,1)$ by Fact \claimref{rule2.imported-rule-2-matrix-units-forgotten} and $\Cl(2,0)$ by Lemma \claimref{rule2.even-subalgebra-one-dimensional-extension}, which is $\Mat_2(\R)$ by \claimref{std-algebra.cl11-matrix}. $e_t = (e_xe_y)^{-1}\omega = -e_xe_y\omega$ since $(e_xe_y)^2 = -1$, so $\varepsilon_\pm e_t = -e_xe_y\varepsilon_\pm\omega = \mp e_xe_y\varepsilon_\pm$. The reflection along $e_x$ is $\hat{\cdot}\circ\Ad_{e_x}$ by the argument of Theorem \claimref{rule2.rule-2-kaluza-klein-algebra}(iii) with $e_x$ in place of $e_4$, and is outer for the same reason. (ii) Fact \claimref{rule2.imported-rule-2-matrix-units-forgotten}, Lemma \claimref{rule2.even-subalgebra-one-dimensional-extension}, and the computations of Theorem \claimref{rule2.rule-2-kaluza-klein-algebra}(v)-(vi) with $\nu^2 = +1$ now the case of Theorem \claimref{rule2.where-obstruction-sits-commuting-square}(c) that it calls spacelike (extension generator of square $-1$, here $e_t$, the time of this plane's labelling): $(\nu e_x)^2 = -\nu^2e_x^2 = -\nu^2$. Theorem \claimref{rule2.where-obstruction-sits-commuting-square}(c) at $(p,q) = (2,0)$, residue $2 \equiv 2 \pmod 4$, whose obstructing extension is the one by a generator of square $-1$, here $e_t$, and whose opposite-sign repair is $\nu^2 = +1$, gives $\Cl(3,1)$ as the Rule 1 ascent of $\Cl(2,0)$. (iii) $\omega^2 = (-1)^3(-1)^2 = -1$ for $\Cl(1,2)$, $p - q = -1 \equiv 3 \pmod 4$; the isomorphism with $\Mat_2(\C)$ is the classification, \claimref{primer.classification-real-clifford-algebras}. The script checks (i)-(iii) in full, the residues included.
